% ----------------------------------------------------------
% RESUMOS
% ----------------------------------------------------------

\setlength{\absparsep}{18pt}

\begin{resumo}
Este trabalho apresenta uma metodologia para avaliar informação visual e inercial de um veículo aéreo não tripulado antes de seu uso em navegação autônoma. A análise reúne dois eixos. No primeiro, 103 execuções completas e 8.305 intervalos sincronizados foram usados para estudar mudanças de proximidade a partir de fluxo óptico, movimento, unidade de medição inercial, comandos e diferença de imagem. Uma perceptron multicamadas foi comparada com referências constantes e modelos de árvores, usando divisão por execução, conjuntos fixos, várias sementes, eventos, gradientes e ablação. No teste com 102 execuções elegíveis, as correlações da rede ficaram entre 0,718 e 0,843; os descritores de diferença de imagem foram o grupo com perda mais consistente quando removido. No segundo eixo, a profundidade foi reprojetada em uma grade tridimensional de ocupação e usada por um planejador A*. Com profundidade de referência do Gazebo, dez de dez runs completaram a rota, com 30 objetivos alcançados e nenhum caminho inseguro adotado. Oito versões de um Depth Anything V2 métrico foram então avaliadas. Embora sete passassem pelo critério de profundidade por pixel, nenhuma passou pelo portão espacial completo por baixa disponibilidade de planos, falso espaço livre ou colisões. Por isso, o voo monocular não foi autorizado. O resultado mostra que métricas por pixel não bastam para liberar um estimador destinado ao planejamento e que a consequência espacial do erro precisa fazer parte da validação.

\vspace{\onelineskip}

\noindent
\textbf{Palavras-chave}: navegação autônoma; visão monocular; profundidade; mapeamento tridimensional; veículo aéreo não tripulado.
\end{resumo}

\begin{resumo}[Abstract]
\begin{otherlanguage*}{english}
This work presents a method for evaluating visual and inertial information from an unmanned aerial vehicle before using it for autonomous navigation. The study has two experimental parts. First, 103 complete runs and 8,305 synchronized intervals were used to analyze proximity changes from optical flow, motion, inertial measurement, commands, and image-difference descriptors. A multilayer perceptron was compared with constant and tree-based baselines using run-level splits, fixed validation and test sets, multiple seeds, event analysis, gradients, and feature-group ablation. With 102 eligible runs, test correlations ranged from 0.718 to 0.843, and removing image-difference descriptors produced the most consistent loss. Second, depth was projected into a three-dimensional occupancy grid and used by an A* planner. With Gazebo reference depth, all ten runs completed the route, reaching 30 global goals without adopting an unsafe path. Eight metric Depth Anything V2 versions were then evaluated. Seven passed the pixel-level depth criterion, but none passed the complete spatial gate because of insufficient path availability, false-free space, or collisions. Monocular flight was therefore not authorized. The results show that pixel metrics alone are insufficient to release a depth estimator for planning and that the spatial consequence of its errors must be included in validation.

\vspace{\onelineskip}

\noindent
\textbf{Keywords}: autonomous navigation; monocular vision; depth estimation; three-dimensional mapping; unmanned aerial vehicle.
\end{otherlanguage*}
\end{resumo}
